\section{Πρόβλημα 6}

\subsection{Ερώτηση (i)}

Θεωρούμε ένα $n \times n$ παίγνιο όπου όλες οι ωφέλειες είναι στο $[0,1]$.
Κατασκευάζουμε το εξής προφίλ μεικτών στρατηγικών: ο παίκτης 1 επιλέγει τις
γραμμές $i$ και $j$ ισοπίθανα, και ο παίκτης 2 επιλέγει τις στήλες $k$ και
$\ell$ ισοπίθανα, με την επιπλέον συνθήκη ότι η $i$ είναι βέλτιστη απόκριση
στη στήλη $k$ και η $\ell$ είναι βέλτιστη απόκριση στη γραμμή $j$.

Συμβολίζουμε με $A$ και $B$ τους πίνακες ωφέλειας των δύο παικτών. Η αναμενόμενη
ωφέλεια των παικτών στο προφίλ $(\p, \q) = \left(\frac{1}{2}e^i +
\frac{1}{2}e^j,\, \frac{1}{2}e^k + \frac{1}{2}e^\ell\right)$ είναι:
$$
u_1(\p, \q) = \frac{1}{4}(a_{ik} + a_{i\ell} + a_{jk} + a_{j\ell}) \\
$$
$$
u_2(\p, \q) = \frac{1}{4}(b_{ik} + b_{i\ell} + b_{jk} + b_{j\ell})
$$

\paragraph{Απόκλιση παίκτη 1.}
Έστω $r$ οποιαδήποτε γραμμή. Η ωφέλεια του παίκτη 1 αν αποκλίνει στη γραμμή $r$
είναι $u_1(e^r, \q) = \frac{1}{2}a_{rk} + \frac{1}{2}a_{r\ell}$. Δηλαδή, το
κέρδος θα είναι
$$
u_1(e^r, \q) - u_1(\p, \q) =
    \frac{1}{2}(a_{rk} + a_{r\ell} -
        \frac{1}{2}a_{ik} - \frac{1}{2}a_{i\ell} -
        \frac{1}{2}a_{jk} - \frac{1}{2}a_{j\ell})
$$
Εφόσον η $i$ είναι βέλτιστη απόκριση στη στήλη $k$, ισχύει $a_{rk} \leq a_{ik}$.
Επίσης όλες οι ώφέλειες είναι στο $[0,1]$, άρα $a_{rk} - \frac{1}{2} a_{ik}
\leq \frac{1}{2} a_{ik} \leq \frac{1}{2}$, $a_{r\ell} \leq 1$ και
$a_{i\ell}, a_{jk}, a_{j\ell} \geq 0$. Συνδυάζοντας
αυτά, έχουμε:
$$
u_1(e^r, \q) - u_1(\p, \q) \leq \frac{1}{2} (\frac{1}{2} + 1 - 0) = \frac{3}{4}
$$

Συμμετρικά, με την ίδια ακριβώς ανάλυση, θα έχουμε το ίδιο για τον παίκτη 2.
Δείξαμε δηλαδή ότι δεν μπορεί να κερδίσει περισσότερο από $\frac{3}{4}$. Για να
δείξουμε ότι το $\frac{3}{4}$ είναι και το μικρότερο φράγμα, αρκεί να βρούμε
ένα παράδειγμα που επιτυγχάνει αυτό το κέρδος. Για
$a_{ik} = a_{rk} = a_{r\ell} = 1$, $a_{i\ell} = a_{jk} = a_{j\ell} = 0$ έχουμε
$u_1(e^r, \q) = \frac{1}{2} + \frac{1}{2} = 1$ και
$u_1(\p, \q) = \frac{1}{4} (1 + 0 + 0 + 0) = \frac{1}{4}$. Άρα,
$\varepsilon = \frac{3}{4}$.

\subsection{Ερώτηση (ii)}

Θεωρούμε ένα $n \times n$ παίγνιο 2 παικτών με $k$-sparse πίνακες $A, B$, όπου
κάθε στοιχείο ανήκει στο διάστημα $[0,1]$. Θα δείξουμε ότι για κάθε
$\varepsilon \geq k/n$, το προφίλ
$(\mathbf{p}, \mathbf{p})$ με $\mathbf{p} = (1/n, \ldots, 1/n)$ είναι
$\varepsilon$-σημείο ισορροπίας.

\paragraph{Απόκλιση παίκτη 1.}
Αρκεί να δείξουμε ότι για κάθε αμιγή στρατηγική $e^r$ του παίκτη 1 ισχύει
$$
    u_1(e^r, \p) - u_1(\p, \p) \leq \varepsilon \,.
$$
Έχουμε
$$
    u_1(\p, \p) = \frac{1}{n^2} \sum_{i=1}^{n} \sum_{j=1}^{n} a_{ij} \geq 0
$$
και
$$
    u_1(e^r, \p) = \frac{1}{n} \sum_{j=1}^{n} a_{rj} \leq \frac{k}{n} \,.
$$

Η απόκλιση θα είναι
$$
    u_1(e^r, \p) - u_1(\p, \p) \leq \frac{k}{n} - 0 = \frac{k}{n} = \varepsilon \,.
$$
